
\documentclass[addpoints]{exam}
\usepackage{amsmath}
\usepackage{amssymb}
\usepackage{color}
\usepackage[version=4]{mhchem}
\usepackage{siunitx}

% \printanswers

\newcommand{\event}{Gases \& Gas Laws}
\newcommand{\tournament}{}
\def\type{0}

\setlength\fillinlinelength{\textwidth}
\CorrectChoiceEmphasis{\color{red}\bfseries}
\SolutionEmphasis{\color{red}\bfseries}

\setlength\answerclearance{2pt}
\pagestyle{headandfoot}
\runningheadrule
\runningheader{\event}{\tournament}{\ifnum\type=1 Class Set Test \else Full Name:\hspace{1cm} \fi}
\runningfootrule
\runningfooter{}{Page \thepage\ of \numpages}{}

\begin{document}
\begin{center}
    {\huge Grade 11 Chemistry \\ \event
    \ifprintanswers \space Key
    \else \space \ifnum\type<2 Test \fi \ifnum\type=1 \textbf{(CLASS SET)} \fi
          \ifnum\type=2 Answer Sheet \fi
    \fi \\ \vspace{3mm}\tournament\vspace{1mm}}
\end{center}

\vspace{.25\textheight}

\begin{center}
    First Name: \ifprintanswers
    \underline{\hspace{3.5cm}}\textbf{KEY}\underline{\hspace{5.5cm}}\\\vspace{5mm}
    \else \underline{\hspace{10cm}}\\ \vspace{5mm} \fi
    Last Name: \underline{\hspace{10cm}}\\ \vspace{5mm}
\end{center}

\noindent\textbf{\underline{Directions:}}
\begin{itemize}
    \ifnum\type=1 \item \textbf{This test is a class set.} Do not write on this paper. \fi
    \item Show all work with units. Use $R = \SI{8.314}{kPa\cdot L/(mol\cdot K)}$.
    STP: $T = \SI{273.15}{K}$, $P = \SI{101.325}{kPa}$.
    \item Convert all temperatures to Kelvin: $T(\text{K}) = T(\si{\celsius}) + 273$.
    \item The test is designed to be completed in 75 minutes.
\end{itemize}
\begin{center}
\textit{For grading use only}\\
\multirowgradetable{1}[pages]
\end{center}

\newpage
\begin{questions}

\section*{Part A --- Multiple Choice (15 marks)}
\vspace{5mm}

\question[1] According to Boyle's Law, at constant temperature, the pressure and volume
of a gas are:
\begin{choices}
    \choice Directly proportional
    \correctchoice Inversely proportional
    \choice Independent of each other
    \choice Equal
\end{choices}

\question[1] A gas occupies $\SI{4.0}{L}$ at $\SI{200}{kPa}$. At constant temperature,
what volume does it occupy at $\SI{100}{kPa}$?
\begin{choices}
    \choice \SI{2.0}{L}
    \correctchoice \SI{8.0}{L}
    \choice \SI{1.0}{L}
    \choice \SI{400}{L}
\end{choices}

\question[1] Charles's Law states that at constant pressure, the volume of a gas is:
\begin{choices}
    \choice Inversely proportional to temperature in Celsius
    \correctchoice Directly proportional to temperature in Kelvin
    \choice Independent of temperature
    \choice Inversely proportional to temperature in Kelvin
\end{choices}

\question[1] A gas at $\SI{27}{\celsius}$ occupies $\SI{3.0}{L}$. At constant pressure, what
volume does it occupy at $\SI{127}{\celsius}$?
\begin{choices}
    \choice \SI{2.25}{L}
    \correctchoice \SI{4.0}{L}
    \choice \SI{3.0}{L}
    \choice \SI{14.1}{L}
\end{choices}

\question[1] Gay-Lussac's Law relates, at constant volume:
\begin{choices}
    \choice Pressure and volume
    \choice Volume and temperature
    \correctchoice Pressure and temperature (in Kelvin)
    \choice Pressure and moles
\end{choices}

\question[1] The combined gas law is:
\begin{choices}
    \choice $PV = nRT$
    \choice $\dfrac{P_1V_1}{T_1} = P_2V_2T_2$
    \correctchoice $\dfrac{P_1V_1}{T_1} = \dfrac{P_2V_2}{T_2}$
    \choice $P_1V_1 = P_2T_2$
\end{choices}

\question[1] The ideal gas law is $PV = nRT$. If $P$ is in kPa, $V$ is in L, and $T$
is in K, then $R =$
\begin{choices}
    \choice \SI{0.08206}{L\cdot atm/(mol\cdot K)}
    \correctchoice \SI{8.314}{kPa\cdot L/(mol\cdot K)}
    \choice \SI{8.314}{J/(g\cdot K)}
    \choice \SI{22.4}{L/mol}
\end{choices}

\question[1] At STP, $\SI{2.0}{mol}$ of any ideal gas occupies:
\begin{choices}
    \choice \SI{22.4}{L}
    \correctchoice \SI{44.8}{L}
    \choice \SI{11.2}{L}
    \choice \SI{24.0}{L}
\end{choices}

\question[1] A gas exerts a pressure of $\SI{80}{kPa}$ at $\SI{27}{\celsius}$.
At constant volume, what pressure does it exert at $\SI{127}{\celsius}$?
\begin{choices}
    \choice \SI{320}{kPa}
    \correctchoice \SI{\approx 106.7}{kPa}
    \choice \SI{60}{kPa}
    \choice \SI{160}{kPa}
\end{choices}

\question[1] Dalton's Law of Partial Pressures states that:
\begin{choices}
    \choice Each gas in a mixture obeys a different gas law
    \correctchoice The total pressure of a gas mixture equals the sum of the partial pressures
    \choice Partial pressures must all be equal
    \choice Lighter gases exert greater pressure
\end{choices}

\question[1] A container holds \ce{N2} at \SI{40}{kPa} and \ce{O2} at \SI{60}{kPa}.
The total pressure is:
\begin{choices}
    \choice \SI{40}{kPa}
    \choice \SI{60}{kPa}
    \correctchoice \SI{100}{kPa}
    \choice \SI{20}{kPa}
\end{choices}

\question[1] Absolute zero ($\SI{0}{K}$) is significant because:
\begin{choices}
    \choice All gases condense to solids
    \correctchoice Theoretically, gas particles have no kinetic energy and volume approaches zero
    \choice Pressure becomes infinite
    \choice It equals $-100\si{\celsius}$
\end{choices}

\question[1] Which variable is held constant in Boyle's Law?
\begin{choices}
    \choice Pressure
    \choice Volume
    \correctchoice Temperature and amount of gas
    \choice Amount of gas only
\end{choices}

\question[1] What is the molar mass of a gas if $\SI{1.50}{g}$ occupies $\SI{0.800}{L}$
at $\SI{101.3}{kPa}$ and $\SI{273}{K}$?
\begin{choices}
    \choice \SI{28.0}{g/mol}
    \correctchoice \SI{42.0}{g/mol}
    \choice \SI{22.4}{g/mol}
    \choice \SI{56.0}{g/mol}
\end{choices}

\question[1] A sealed syringe contains $\SI{20.0}{mL}$ of gas at $\SI{100}{kPa}$.
The plunger is pushed until the pressure is $\SI{250}{kPa}$. The new volume is:
\begin{choices}
    \choice \SI{50.0}{mL}
    \correctchoice \SI{8.00}{mL}
    \choice \SI{25.0}{mL}
    \choice \SI{10.0}{mL}
\end{choices}


\section*{Part B --- Short Answer (35 marks)}

\question Apply the appropriate gas law to each problem.
\begin{parts}
    \part[3] \textbf{(Boyle's Law)} A scuba diver's lungs hold $\SI{6.0}{L}$ of air at a
    depth where the pressure is $\SI{300}{kPa}$. If she ascends to the surface where the
    pressure is $\SI{101.3}{kPa}$ without exhaling, what volume would her lungs reach?
    (This is why divers must exhale while ascending.)
    \part[3] \textbf{(Charles's Law)} A balloon has a volume of $\SI{2.50}{L}$ at
    $\SI{20.0}{\celsius}$. It is placed in a freezer at $\SI{-15.0}{\celsius}$.
    What is the new volume at constant pressure?
    \part[3] \textbf{(Gay-Lussac's Law)} A car tire contains air at $\SI{220}{kPa}$ and
    $\SI{15}{\celsius}$. After driving, the tire temperature rises to $\SI{50}{\celsius}$.
    What is the new pressure? (Assume constant volume.)
\end{parts}
\begin{solution}[10cm]
\textbf{(a)} $T$ constant. $P_1V_1 = P_2V_2$:
\[
V_2 = \frac{P_1V_1}{P_2} = \frac{300 \times 6.0}{101.3} = \mathbf{\SI{17.8}{L}}
\]
This dangerously exceeds lung capacity, illustrating why divers must exhale while ascending.

\textbf{(b)} $T_1 = 293\,\text{K}$,\; $T_2 = 258\,\text{K}$.
\[
V_2 = V_1 \cdot \frac{T_2}{T_1} = 2.50 \times \frac{258}{293} = \mathbf{\SI{2.20}{L}}
\]

\textbf{(c)} $T_1 = 288\,\text{K}$,\; $T_2 = 323\,\text{K}$.
\[
P_2 = P_1 \cdot \frac{T_2}{T_1} = 220 \times \frac{323}{288} = \mathbf{\SI{247}{kPa}}
\]
\end{solution}

\question A sample of gas occupies $\SI{1.50}{L}$ at $\SI{25}{\celsius}$ and
$\SI{150}{kPa}$.
\begin{parts}
    \part[3] Use the combined gas law to find the volume at $\SI{0}{\celsius}$ and
    $\SI{100}{kPa}$.
    \part[3] Use the ideal gas law to find the number of moles of gas in the original
    sample ($T = \SI{25}{\celsius}$, $P = \SI{150}{kPa}$, $V = \SI{1.50}{L}$).
    \part[2] Calculate the molar mass if the sample has a mass of $\SI{4.80}{g}$.
\end{parts}
\begin{solution}[9cm]
\textbf{(a)} $T_1 = 298\,\text{K}$,\; $T_2 = 273\,\text{K}$.
\[
V_2 = V_1 \cdot \frac{P_1}{P_2} \cdot \frac{T_2}{T_1}
= 1.50 \times \frac{150}{100} \times \frac{273}{298} = \mathbf{\SI{2.06}{L}}
\]

\textbf{(b)}
\[
n = \frac{PV}{RT} = \frac{150 \times 1.50}{8.314 \times 298} = \frac{225}{2477.6}
= \mathbf{0.0908 \text{ mol}}
\]

\textbf{(c)}
\[
M = \frac{m}{n} = \frac{4.80}{0.0908} = \mathbf{\SI{52.9}{g/mol}}
\]
\end{solution}

\question \textbf{Gas stoichiometry.} Zinc reacts with hydrochloric acid:
\[\ce{Zn(s) + 2HCl(aq) -> ZnCl2(aq) + H2(g)}\]
\begin{parts}
    \part[3] What volume of \ce{H2} gas is produced at STP when $\SI{6.54}{g}$ of Zn
    ($M = \SI{65.4}{g/mol}$) reacts completely?
    \part[3] At what temperature (in \si{\celsius}) would this volume of \ce{H2} occupy
    $\SI{3.00}{L}$ at $\SI{101.3}{kPa}$?
\end{parts}
\begin{solution}[9cm]
\textbf{(a)} $n(\ce{Zn}) = \dfrac{6.54}{65.4} = 0.100$ mol;\quad
$n(\ce{H2}) = 0.100$ mol (1:1 ratio)

At STP: $V = n \times \SI{22.4}{L/mol} = 0.100 \times 22.4 = \mathbf{\SI{2.24}{L}}$

\textbf{(b)} Using ideal gas law with $n = 0.100$ mol:
\[
T = \frac{PV}{nR} = \frac{101.3 \times 3.00}{0.100 \times 8.314}
= \frac{303.9}{0.8314} = \SI{365.5}{K}
\]
\[
T = 365.5 - 273 = \mathbf{\SI{92.5}{\celsius}}
\]
\end{solution}

\question \textbf{Dalton's Law and moist gas.} A gas is collected over water at
$\SI{25}{\celsius}$ and a total pressure of $\SI{101.3}{kPa}$. The vapour pressure of
water at $\SI{25}{\celsius}$ is $\SI{3.17}{kPa}$.
\begin{parts}
    \part[2] Calculate the partial pressure of the dry gas collected.
    \part[3] If $\SI{0.450}{L}$ of the moist gas is collected, calculate the moles of
    \textbf{dry gas} present ($T = \SI{25}{\celsius} = \SI{298}{K}$).
\end{parts}
\begin{solution}[7cm]
\textbf{(a)} By Dalton's Law:
\[
P_\text{gas} = P_\text{total} - P_{\ce{H2O}} = 101.3 - 3.17 = \mathbf{\SI{98.1}{kPa}}
\]

\textbf{(b)}
\[
n = \frac{P_\text{gas} \cdot V}{RT} = \frac{98.1 \times 0.450}{8.314 \times 298}
= \frac{44.15}{2477.6} = \mathbf{0.01782 \text{ mol}}
\]
\end{solution}

\end{questions}
\end{document}
